% Section 9 — Discussion
\section{Discussion}
\label{sec:discussion}

\paragraph{Two routes to the same artifact} The pilots suggest a picture different from the one we
expected when we began. Frontier models, reasoning only in language and never seeing a dataset, write
regression algorithms that work, and the best of them ranked first among eighteen methods on our
small panel. A data-conditioned engine writes a different algorithm for each dataset and, with a few
hundred evaluations, is about as good. These are two routes to the same kind of artifact --- a
standalone, readable program --- with opposite economics. Route~1 pays a large cost once, in
language, and produces a general algorithm; route~2 pays per dataset, in program evaluations, and
produces a specific one. Neither result is a track-C result, and neither settles our question. What
they establish is that the question is the right one: both routes already produce strong programs,
and what remains unknown is whether the second can propose its program before testing anything.

\paragraph{The wrong objective, stated carefully} The claim that language models ``chase the wrong
objective'' is usually made too broadly. The careful version is \cref{tab:objectives}. A pretrained
model maximizes the likelihood of text people wrote; a preference-tuned model maximizes a model of
what people prefer \citep{christiano2017rlhf,ouyang2022instructgpt}; a reasoning-tuned model is also
rewarded by checkers on problems whose answers can be verified \citep{openai2024o1,deepseek2025r1},
and today's frontier models combine all three. None of these objectives refers to what a program
does on data no one has seen, and the recipes behind the four entrants are not documented, so we
infer nothing about them from their names. Sutton's version is that language models have no ground truth
\citep{sutton2025dwarkesh}; Almeida's, that preference is the wrong target for systems that act on
their own \citep{almeida2026aie}; LeCun's, that the architecture itself is the problem
\citep{lecun2026mittr}. Our evidence is consistent with a narrower reading than any of these. Frontier
models are not unable to produce working new algorithms; they produce the algorithm the field would
write. Training on likelihood and preference draws them to the centre of what is known
\citep{grinsztajn2022trees}, and the four entrants' convergence on one design is that pull, observed.
An objective defined by measured outcome (\cref{eq:j1}) pulls toward whatever works, including what
nobody wrote. Regression is a domain where that objective is cheap to compute --- run the program,
measure the error --- and therefore one where it can be tested today.

\paragraph{The Eureka Machine and amortized discovery} Socher's \emph{Eureka Machine}
\citep{socher2026eureka} rests discovery on four pillars --- a map of human knowledge, a model of
physical reality learned from measurement, high-fidelity digital twins and autonomous laboratories
--- above which, in his words, ``millions of AI agents will coordinate like a scientific community''
\citep{socher2026bookbite}. Recursive describes parallel research threads that propose, implement,
run and validate experiments \citep{recursive2026firststeps}, and reports using its automated
research system of the same name to build its Flint model \citep{recursive2026flint}; these are
company reports of AI research and development, not independent demonstrations of the book's
measurement, simulation and laboratory pillars, nor of unrestricted self-improvement. SwarmRSI
likewise lets experimental feedback steer a swarm of scientific agents \citep{meng2026swarmrsi}. InnovationZero is
complementary and, on one axis, opposite. It takes the second pillar --- measurement data --- as its
input, and asks whether the \emph{output} of the loop, a method that works, can be produced in one
pass for problems like those it has already seen. The loop is how a field discovers; amortization is
how it stops paying for discovery it has already made. The two compose: an iterative research
system produces exactly the measured history an amortized proposer learns from, and a better first
proposal gives the loop a better place to start. The AlphaGo lineage did both: search discovered,
the network remembered.

\paragraph{The cost of invention, counted honestly} Token yield (\cref{sec:yield}) is often read as an
efficiency score. It is better read as a map of where a system explores. The frontier sessions
explored in text and kept a tenth of what they processed; the engine explored in program space and
generated no text until it committed. That is a structural difference, not yet an economic one. A
frontier model's authoring cost buys a program that serves every later dataset, and on our pilot
those programs fitted in tenths of a second; the engine paid seconds per dataset. Whether learned
first proposals are cheaper per useful algorithm once training is amortized is gate M6, and we have
not measured it. Cheaper invention would also not mean less of it: Jevons observed in 1865 that more
efficient coal engines increased the consumption of coal \citep{jevons1865}, and the same is likely
to be true of algorithms.

\paragraph{What each outcome would mean} The protocol is built so that every outcome is
informative. If the learned proposer beats selection and retrieval over the same archive, that
establishes a predictive advantage under that comparison; only archive exclusion and mechanism tests
can then attribute it to construction rather than recall. If a distilled bank or retrieval matches
it, the evidence supports amortized selection, which is useful but not invention. If a few target
evaluations erase its advantage, the advantage belongs to the low-feedback regime alone. None of
these results would establish new mathematics, general intelligence or scientific autonomy, and
each answers the paper's question more clearly than a leaderboard rank.

\paragraph{Limitations} The operations from which kernels are built are human-designed: ridge
regression, Gaussian processes and gradient boosting are decades of statistics, not rules of a game.
``New'' in this paper therefore means a new composition for a dataset, graded by the ladder in
\cref{tab:novelty}, not new mathematics. Only regression is studied. The frontier entrants are single,
non-pristine attempts, and every token count is an estimate. The pilots are small: ten datasets in
one case, and in the other 64 training families with prepared inner-validation contexts, preprocessing
fitted before the original split, and three seeds fixed after the first had been inspected. The
engine's timings and the frontier entrants' come from different hosts. Most of the training corpus
has unknown licence status and cannot be redistributed; the replication panels are public datasets.
Elo is relative to the panel and the pool of opponents, and we report magnitudes beside it. Finally,
the result this paper is about has not been run: no neural proposer in this system has yet met the
track-C contract on held-out data.
